\section{Conclusion}
\label{sec:conclusion}

We asked, in a preregistered experiment, whether language models selectively erode their commitments to human oversight when they repeatedly revise the written principles that govern them.
\begin{itemize}
\item \textbf{In our setting, they did not.} Oversight commitments eroded no faster than matched restraints on the models' own agency, and probably more slowly, although this depends on particular items.
\item \textbf{Erosion concentrated in restraints on self-expression,} including political neutrality.
\item \textbf{Human coders judged many of those changes to be improvements,} but judged weakened oversight commitments to be worse.
\item \textbf{The oversight commitment eroded most often}, accepting changes to one's own values, was typically made conditional on the model's own safety principles, and human coders disagreed about whether that is good.
\end{itemize}

Written principles are becoming a central instrument of AI governance. Our results suggest that, as models help to revise these documents, it will be important to think carefully about how to manage AI-assisted changes to written principles, including through record keeping, human review and approval, and special attention to changes that shift decision-making authority away from people and towards AI systems. The preregistration, code, generations, judge codes and human codes are publicly available to support future research.
